% Complete typesetting of the unchanged public Markdown proof.
% Original: https://github.com/KokunoYumeto/zeta-function-research-reader/blob/97bc5b1a25edab0d3fcc8c43b79cfa8ece8ce76c/workbenches/splitzero-tandem/branches/tau-zero-prime-positivity/supporting_proofs/TAU_PRIME_SPECTRUM_DERIVATION.md
% Source SHA256: b0bca40fdfbc1d49cd8620e6987e53e810d6c86b413001a28075011f974e46b8
% No mathematical content or citations were added or removed.
% Options for packages loaded elsewhere
\PassOptionsToPackage{unicode}{hyperref}
\PassOptionsToPackage{hyphens}{url}
\documentclass[
  10pt,
]{article}
\usepackage{xcolor}
\usepackage[margin=22mm]{geometry}
\usepackage{amsmath,amssymb}
\setcounter{secnumdepth}{-\maxdimen} % remove section numbering
\usepackage{iftex}
\ifPDFTeX
  \usepackage[T1]{fontenc}
  \usepackage[utf8]{inputenc}
  \usepackage{textcomp} % provide euro and other symbols
\else % if luatex or xetex
  \usepackage{unicode-math} % this also loads fontspec
  \defaultfontfeatures{Scale=MatchLowercase}
  \defaultfontfeatures[\rmfamily]{Ligatures=TeX,Scale=1}
\fi
\usepackage{lmodern}
\ifPDFTeX\else
  % xetex/luatex font selection
\fi
% Use upquote if available, for straight quotes in verbatim environments
\IfFileExists{upquote.sty}{\usepackage{upquote}}{}
\IfFileExists{microtype.sty}{% use microtype if available
  \usepackage[]{microtype}
  \UseMicrotypeSet[protrusion]{basicmath} % disable protrusion for tt fonts
}{}
\makeatletter
\@ifundefined{KOMAClassName}{% if non-KOMA class
  \IfFileExists{parskip.sty}{%
    \usepackage{parskip}
  }{% else
    \setlength{\parindent}{0pt}
    \setlength{\parskip}{6pt plus 2pt minus 1pt}}
}{% if KOMA class
  \KOMAoptions{parskip=half}}
\makeatother
\usepackage{longtable,booktabs,array}
\newcounter{none} % for unnumbered tables
\usepackage{calc} % for calculating minipage widths
% Correct order of tables after \paragraph or \subparagraph
\usepackage{etoolbox}
\makeatletter
\patchcmd\longtable{\par}{\if@noskipsec\mbox{}\fi\par}{}{}
\makeatother
% Allow footnotes in longtable head/foot
\IfFileExists{footnotehyper.sty}{\usepackage{footnotehyper}}{\usepackage{footnote}}
\makesavenoteenv{longtable}
\setlength{\emergencystretch}{3em} % prevent overfull lines
\providecommand{\tightlist}{%
  \setlength{\itemsep}{0pt}\setlength{\parskip}{0pt}}
\usepackage{mathrsfs}
\usepackage{mathtools}
\usepackage{xurl}
\emergencystretch=3em
\setcounter{tocdepth}{1}
\newcommand{\Beta}{\mathrm{B}}
\DeclareUnicodeCharacter{00B2}{\ensuremath{^2}}
\DeclareUnicodeCharacter{00B1}{\ensuremath{\pm}}
\DeclareUnicodeCharacter{00D7}{\ensuremath{\times}}
\DeclareUnicodeCharacter{2192}{\ensuremath{\to}}
\DeclareUnicodeCharacter{21A6}{\ensuremath{\mapsto}}
\DeclareUnicodeCharacter{2208}{\ensuremath{\in}}
\DeclareUnicodeCharacter{2209}{\ensuremath{\notin}}
\DeclareUnicodeCharacter{2212}{\ensuremath{-}}
\DeclareUnicodeCharacter{2227}{\ensuremath{\wedge}}
\DeclareUnicodeCharacter{2228}{\ensuremath{\vee}}
\DeclareUnicodeCharacter{2260}{\ensuremath{\ne}}
\DeclareUnicodeCharacter{2264}{\ensuremath{\le}}
\DeclareUnicodeCharacter{2265}{\ensuremath{\ge}}
\DeclareUnicodeCharacter{2297}{\ensuremath{\otimes}}
\usepackage{bookmark}
\IfFileExists{xurl.sty}{\usepackage{xurl}}{} % add URL line breaks if available
\urlstyle{same}
\hypersetup{
  hidelinks,
  pdfcreator={LaTeX via pandoc}}

\author{}
\date{}

\begin{document}

\section{The original tau spectrum, its residues and its prime-counting
receivers}\label{the-original-tau-spectrum-its-residues-and-its-prime-counting-receivers}

22 September 2026. TPS1--TPS59. This note starts with the original
semiring and retains its arithmetic and support coordinates. The ideal
spectrum, congruence quotients, stalks, monoid spectrum and contracted
monoid-ring splitting are existing programme results, verified below
rather than presented as new classifications. The residue comparisons,
their exact support pullbacks, and the analytic correction for the
explicitly defined supported prime-residue product are developed with
complete proofs.

Source coverage:
\href{https://github.com/KokunoYumeto/zeta-function-research-reader/blob/ac168d55cbb400a6c8926a00d67909ba65646ac5/workbenches/splitzero-tandem/branches/identity-absorber-square/sources/17555345_11.tex}{the
original programme paper, 11.tex}, construction\_S, def:addition\_S,
def:multiplication\_S, thm:S\_is\_semiring and prime classification,
source lines578--1005; and \href{globalization_note.tex}{A Note on the
Globalization Semiring G(Z), complete source}, ideal, congruence,
localization, stalk, monoid-spectrum and base-extension proofs. These
earlier sources retain their identity. The mathematical derivation below
includes every connecting calculation it uses.

\subsection{1. Original operations and both coordinate
maps}\label{original-operations-and-both-coordinate-maps}

Let \[
S=G(\mathbb Z)=\{\tau\}\sqcup\mathbb Z^\bullet,
\qquad e=0^\bullet,\qquad 1_S=1^\bullet.
\tag{TPS1}
\] The superscript records a supported integer; \(\tau\) is external,
and in particular \(\tau\ne e\). The original operations are \[
m^\bullet+n^\bullet=(m+n)^\bullet,\qquad
m^\bullet n^\bullet=(mn)^\bullet,\qquad
x+\tau=x,\qquad x\tau=\tau.
\tag{TPS2}
\] Define \(\mathbb B=\{0,1\}\) with OR as addition and AND as
multiplication. The two unital semiring maps are \[
\pi:S\to\mathbb Z,\quad \pi(\tau)=0,\quad\pi(n^\bullet)=n,
\qquad
\chi:S\to\mathbb B,\quad \chi(\tau)=0,\quad\chi(n^\bullet)=1.
\tag{TPS3}
\] Indeed \((\pi,\chi)\) identifies \(S\) with
\(\{(0,0)\}\cup(\mathbb Z\times\{1\})\subset\mathbb Z\times\mathbb B\).
The sum and product of two supported pairs are \((m+n,1)\) and
\((mn,1)\); adding \((0,0)\) fixes either pair, and multiplying by it
gives \((0,0)\). This proves closure, preservation of both operations,
both identities, and injectivity. The semiring axioms follow from those
of the product. For \(\mathbb B\), distributivity follows by checking
whether the factor is \(0\) or \(1\); its other axioms are the
corresponding Boolean operations. In particular \[
\pi^{-1}(0)=\{\tau,e\},\qquad
\chi^{-1}(0)=\{\tau\},\qquad
1^\bullet+(-1)^\bullet=e\ne\tau.
\tag{TPS4}
\] The original element named \(\mathcal T\) is identified with \(\tau\)
by the map fixing every integer, not by the arithmetic quotient \(\pi\).

\subsection{2. All semiring ideals, primes and their
topology}\label{all-semiring-ideals-primes-and-their-topology}

An ideal contains \(\tau\), is closed under addition, and under
multiplication by every element of \(S\). Write \[
I_n=\{\tau\}\cup(n\mathbb Z)^\bullet\quad(n\ge0),
\qquad P_\tau=\{\tau\},\quad P_0=I_0,\quad P_p=I_p
\quad(p\text{ a positive rational prime}).
\tag{TPS5}
\] Every ideal is exactly \(P_\tau\) or exactly one \(I_n\). To prove
it, intersect the ideal with \(\mathbb Z^\bullet\). If this is empty the
ideal is \(P_\tau\). Otherwise multiplication by \(e\) includes \(e\),
multiplication by \((-1)^\bullet\) includes additive inverses in the
supported group, and addition and scalar multiplication make the integer
intersection a ring ideal. If it contains a nonzero integer, take its
least positive integer \(n\); division with remainder proves every
member is divisible by \(n\). If it contains no nonzero integer it is
\(0\mathbb Z\). Conversely each displayed set is closed under all
required operations. These arguments also prove uniqueness and that
\(I_n\) is generated by \(n^\bullet\).

A proper ideal is prime when membership of a product forces membership
of a factor. A product equals \(\tau\) exactly when one factor is
\(\tau\), so \(P_\tau\) is prime. For \(I_n\), products of two supported
elements reduce the test exactly to primeness of \(n\mathbb Z\). If
\(n\ge2\) is composite, a factorization \(n=ab\) with \(1<a,b<n\)
disproves primeness; a prime \(p\) satisfies the test by integer
divisibility. The zero ideal satisfies the test because \(\mathbb Z\) is
a domain; \(I_1=S\) is improper. Thus the full spectrum is \[
\boxed{X=\operatorname{Spec}S=\{P_\tau,P_0\}\sqcup\{P_p:p\text{ prime}\},
\qquad P_\tau\subsetneq P_0\subsetneq P_p.}
\tag{TPS6}
\] Distinct \(P_p\)'s are incomparable. The maximal ideals and closed
points are precisely the \(P_p\)'s, and the longest prime chains have
length two. These recover the existing programme classification.

For the Zariski topology, \(V(I)=\{P:I\subseteq P\}\) and
\(D(a)=X\setminus V((a))\). Direct evaluation gives \[
D(\tau)=\varnothing,\quad D(e)=\{P_\tau\},\quad
D(n^\bullet)=\{P_\tau,P_0\}\cup\{P_p:p\nmid n\}quad(n\ne0).
\tag{TPS7}
\] The closed sets are \(\varnothing\), \(X\),
\(C=\{P_0\}\cup\{P_p\}=V(e)\), and finite sets of closed rational-prime
points. Indeed \(V(P_\tau)=X\), \(V(I_0)=C\), \(V(I_1)=\varnothing\),
and for \(n\ge2\), \(V(I_n)\) consists of the primes dividing \(n\).
Every finite set is obtained from the product of its primes. Thus
\(P_\tau\) is a dense generic point, but also has the open singleton
\(D(e)\); \(P_0\) is generic only for the closed arithmetic subspace
\(C\).

Contraction of ideals under TPS3 gives the complete continuous maps \[
\pi^*:\operatorname{Spec}\mathbb Z\xrightarrow{\sim}C\hookrightarrow X,
\quad (0)\mapsto P_0,\quad(p)\mapsto P_p,
\] \[
\chi^*:\operatorname{Spec}\mathbb B\xrightarrow{\sim}\{P_\tau\}\hookrightarrow X.
\tag{TPS8}
\] The only proper ideal of \(\mathbb B\) is \(\{0\}\), and it is prime.
Preimages of the basic opens are the basic opens of the images of their
defining elements, proving continuity. TPS7 proves that the first map is
a homeomorphism onto the closed subspace and the second onto the open
singleton. Their images are disjoint and exhaust \(X\), but \(X\) is
irreducible because every nonempty open contains \(P_\tau\).

\subsection{3. Localizations, sheaf sections and exact stalk
maps}\label{localizations-sheaf-sections-and-exact-stalk-maps}

For a multiplicative denominator set \(U\subset S\) containing \(1\),
define fractions by \[
(x,u)\sim(y,v)\iff\exists w\in U:\ wvx=wuy.
\tag{TPS9}
\] The fraction operations are \(x/u+y/v=(xv+yu)/(uv)\) and
\((x/u)(y/v)=xy/(uv)\). Multiplying witnesses proves that they respect
this equivalence. If a unital semiring map sends every denominator to a
unit, \(x/u\mapsto f(x)f(u)^{-1}\) is the unique extension: a witness
becomes an equality after its invertible factors are cancelled. This
proves the localization universal property.

If \(U\subset\mathbb Z^\bullet\) avoids \(e\), then all its denominators
and witnesses are nonzero integers. No fraction with supported numerator
equals one with numerator \(\tau\), because every product of supported
integers is supported. For two supported fractions, cancellation of the
nonzero witness is exactly the ordinary integer fraction relation.
Consequently \[
S_U\xrightarrow{\sim}G(\mathbb Z_U),\qquad
\tau/u\mapsto\tau,\quad n^\bullet/u^\bullet\mapsto(n/u)^\bullet.
\tag{TPS10}
\] The fraction operations prove that this bijection is an isomorphism.
If instead \(e\in U\) and \(\tau\notin U\), every two supported
fractions are equal, using witness \(e\) in TPS9; fractions with
numerator \(\tau\) remain distinct from them. The resulting two-element
semiring is \(\mathbb B\), with supported fractions mapping to \(1\). If
\(\tau\in U\), all fractions coincide because the absorbing element has
become a unit, giving the one-element semiring.

In particular the original stalks and their nonunits are

{\def\LTcaptype{none} % do not increment counter
\begin{longtable}[]{@{}
  >{\raggedright\arraybackslash}p{(\linewidth - 6\tabcolsep) * \real{0.2500}}
  >{\raggedright\arraybackslash}p{(\linewidth - 6\tabcolsep) * \real{0.2500}}
  >{\raggedright\arraybackslash}p{(\linewidth - 6\tabcolsep) * \real{0.2500}}
  >{\raggedright\arraybackslash}p{(\linewidth - 6\tabcolsep) * \real{0.2500}}@{}}
\toprule\noalign{}
\begin{minipage}[b]{\linewidth}\raggedright
Point
\end{minipage} & \begin{minipage}[b]{\linewidth}\raggedright
Stalk \(S_P\)
\end{minipage} & \begin{minipage}[b]{\linewidth}\raggedright
Unique maximal ideal
\end{minipage} & \begin{minipage}[b]{\linewidth}\raggedright
Quotient by that ideal
\end{minipage} \\
\midrule\noalign{}
\endhead
\bottomrule\noalign{}
\endlastfoot
\(P_\tau\) & \(\mathbb B\) & \(\{0\}\) & \(\mathbb B\) \\
\(P_0\) & \(G(\mathbb Q)\) & \(\{\tau,0^\bullet\}\) & \(\mathbb Q\) \\
\(P_p\) & \(G(\mathbb Z_{(p)})\) &
\(\{\tau\}\cup(p\mathbb Z_{(p)})^\bullet\) & \(\mathbb F_p\) \\
\end{longtable}
}

\[
\mathbb Z_{(p)}=\{a/b\in\mathbb Q:p\nmid b\}.
\tag{TPS11}
\] The stalk is localization at \(S\setminus P\); TPS10 and its
\(e\)-denominator case give every entry. A supported element of \(G(A)\)
is a unit precisely when its ring value is a unit of \(A\), and \(\tau\)
is never a unit in these nontrivial semirings. The listed nonunits are
ideals and their complements are the units, which proves uniqueness of
the maximal ideal. The last column uses the congruence quotient proved
in the next section.

For completeness, the localization sheaf is explicit on every open. A
nonempty open is \(\{P_\tau\}\), or \(D(n^\bullet)\) for some nonzero
integer \(n\). Its sections are respectively \[
\mathcal O_X(\{P_\tau\})=\mathbb B,\qquad
\mathcal O_X(D(n^\bullet))=G(\mathbb Z[1/n]).
\tag{TPS12}
\] The restrictions between arithmetic opens are the corresponding
rational inclusions; the restriction to \(\{P_\tau\}\) is the support
map, sending every supported fraction to \(1\). The empty open has the
one-element section semiring. These prescriptions form a sheaf. Indeed
any two nonempty opens intersect at \(P_\tau\), so compatible sections
have the same support label. All absent sections glue uniquely to
absence. If they are supported and the union contains \(P_0\), the opens
containing \(P_0\) have rational values that coincide on their
intersections, giving one rational number. For each rational prime
retained in the union, some covering open contains that point, so this
rational number has no denominator factor at that prime. It belongs to
the stated localization of the union. The section on a cover member
\(\{P_\tau\}\) then agrees automatically. A union equal to
\(\{P_\tau\}\) has the two-element sheaf directly. This proves existence
and uniqueness of all gluings and the stated stalk descriptions without
changing their denominators.

All nonidentity generization maps are \[
G(\mathbb Z_{(p)})\hookrightarrow G(\mathbb Q),\qquad
G(\mathbb Q)\xrightarrow{\chi}\mathbb B,
\qquad
G(\mathbb Z_{(p)})\xrightarrow{\chi}\mathbb B.
\tag{TPS13}
\] They follow from inclusion of denominator sets; the last is the
composite of the first two. The morphism \(\pi^*\) has stalk maps
\(G(\mathbb Z_{(p)})\to\mathbb Z_{(p)}\) and
\(G(\mathbb Q)\to\mathbb Q\), which are exactly arithmetic projection.
The stalk map of \(\chi^*\) is \(\mathbb B\to\mathbb B\), the identity.
Thus support and arithmetic are connected by explicitly computed local
maps.

\subsection{4. Both residue congruences and their different zero
fibres}\label{both-residue-congruences-and-their-different-zero-fibres}

For an ideal \(I\), define its Bourne congruence by \[
x\equiv_I y\iff\exists a,b\in I:\ x+a=y+b.
\tag{TPS14}
\] Reflexivity uses \(a=b=\tau\); symmetry interchanges witnesses. If
\(x+a=y+b\) and \(y+c=z+d\), adding gives \(x+(a+c)=z+(b+d)\), proving
transitivity. Adding equalities and multiplying an equality by any
scalar prove compatibility with the semiring operations. It is the
smallest congruence identifying every element of \(I\) with \(\tau\):
the displayed equality proves necessity, while \(a+\tau=\tau+a\) proves
that each \(a\in I\) is identified with \(\tau\).

Define, for \(n\ge0\), with modulus zero meaning equality of integers,
\[
\rho_n:\quad \tau\text{ is a singleton class, and }a^\bullet\sim b^\bullet
\iff a\equiv b\pmod n,
\] \[
\rho'_n:\quad x\sim y\iff\pi(x)\equiv\pi(y)\pmod n.
\tag{TPS15}
\] These are exactly the two congruences over each integer modulus, and
\[
S/\rho_n\cong G(\mathbb Z/n\mathbb Z),\qquad
S/\rho'_n\cong\mathbb Z/n\mathbb Z,
\qquad \equiv_{I_n}=\rho'_n.
\tag{TPS16}
\] To prove completeness, restrict any congruence to the supported
abelian group. The class of \(e\) is an integer ideal, and adding
\((-b)^\bullet\) shows that \(a^\bullet\sim b^\bullet\) exactly when
\((a-b)^\bullet\sim e\). Thus it is reduction modulo a unique \(n\). If
\(\tau\) is equivalent to any supported element, multiplication by \(e\)
gives \(\tau\sim e\); its entire class is then
\(\{\tau\}\cup(n\mathbb Z)^\bullet\). Otherwise its class is a
singleton. These are TPS15, whose quotient operations give the two
stated isomorphisms. For the last equality, reduction modulo \(n\) sends
the witnesses of TPS14 to zero, proving one inclusion. Conversely, if
supported \(a,b\) are congruent, use witnesses \((b-a)^\bullet,e\). If
\(\tau\) and \(b^\bullet\) are congruent, then \(b^\bullet\in I_n\) and
TPS14 identifies it with \(\tau\). This exhausts the pairs.

At a closed prime \(P_p\), the two natural residue maps and their
kernels are consequently \[
q_p^{\mathrm{sup}}:S\twoheadrightarrow G(\mathbb F_p),
\quad \tau\mapsto\tau,\quad n^\bullet\mapsto\bar n^\bullet,
\quad (q_p^{\mathrm{sup}})^{-1}(\tau)=\{\tau\},
\] \[
(q_p^{\mathrm{sup}})^{-1}(0^\bullet)=(p\mathbb Z)^\bullet,
\qquad
q_p^{\mathrm{ar}}=\pi_{\mathbb F_p}q_p^{\mathrm{sup}}:S\twoheadrightarrow\mathbb F_p,
\quad (q_p^{\mathrm{ar}})^{-1}(0)=P_p.
\tag{TPS17}
\] The first quotient has \(p+1\) elements; the second has \(p\). Both
extend to \(G(\mathbb Z_{(p)})\), because denominators prime to \(p\)
map to units. There they send \((a/b)^\bullet\) to
\((\bar a/\bar b)^\bullet\) and \(\bar a/\bar b\), respectively. This
proves all residue assertions in TPS11. The supported residue has the
injective joint coordinate map \[
G(\mathbb F_p)\hookrightarrow\mathbb F_p\times\mathbb B,
\quad \tau\mapsto(0,0),\quad a^\bullet\mapsto(a,1).
\tag{TPS18}
\] The prime map from its two-point semiring spectrum sends \(\{\tau\}\)
to \(P_\tau\) and \(\{\tau,0^\bullet\}\) to \(P_p\). Thus the
support-preserving residue retains an actual generic support point over
the original generic support point. It is not the quotient by the ideal
\(P_p\).

\subsection{5. Composite residues and the exact mixed-support
defect}\label{composite-residues-and-the-exact-mixed-support-defect}

For any two commutative rings \(A,B\), including the zero ring, their
support maps give a canonical isomorphism \[
\boxed{G(A\times B)\xrightarrow{\sim}G(A)\times_{\mathbb B}G(B),
\quad \tau\mapsto(\tau,\tau),\quad(a,b)^\bullet\mapsto(a^\bullet,b^\bullet).}
\tag{TPS19}
\] Equality of the two support characters forces a pair to be either
simultaneously absent or simultaneously supported. These are exactly the
displayed image and its inverse. Both operations and units are preserved
by direct component evaluation. A pair of semiring maps into
\(G(A),G(B)\) with equal support composites has a unique map into this
image, proving the fibre-product universal property.

For coprime positive integers \(m,n\), the integer Chinese remainder map
gives \[
\boxed{G(\mathbb Z/mn\mathbb Z)
\cong G(\mathbb Z/m\mathbb Z)\times_{\mathbb B}G(\mathbb Z/n\mathbb Z).}
\tag{TPS20}
\] For a complete proof of the ring map, choose integers \(u,v\) with
\(um+vn=1\). The residue pair \((a\bmod m,b\bmod n)\) is the image of
\(a v n+b u m\bmod mn\). If both residues of an integer vanish,
coprimality implies divisibility by \(mn\), proving injectivity. The map
preserves addition and multiplication, so it is a ring isomorphism;
TPS19 then proves TPS20.

For finite rings of cardinalities \(m,n\), the image in the full
Cartesian product excludes exactly \[
A^\bullet\times\{\tau\},\qquad \{\tau\}\times B^\bullet,
\] \[
|G(A\times B)|=mn+1,\qquad
|G(A)\times G(B)|=(m+1)(n+1),
\qquad (m+1)(n+1)-(mn+1)=m+n.
\tag{TPS21}
\] These are actual disjoint mixed-support sets of sizes \(m,n\). More
generally, a finite set over \(\mathbb B\) has the two-component count
vector \((|X_0|,|X_1|)\), and fibre products multiply these vectors
componentwise. For \(G(A)\) this vector is \((1,m)\); hence
\((1,m)(1,n)=(1,mn)\). Forgetting the label takes the sum of coordinates
and gives \(m+1\); this summation does not preserve componentwise
multiplication. This identifies the precise map responsible for the
cardinality defect.

\subsection{6. Three specified counting
functions}\label{three-specified-counting-functions}

Ideal multiplication in \(S\) satisfies \(I_mI_n=I_{mn}\): every sum of
products is divisible by \(mn\), and every supported multiple of \(mn\)
is a product of an element of \(I_m\) and one of \(I_n\); \(\tau\) is
included. The Bourne norm of a finite-norm ideal is \[
N_{\mathrm{ar}}(I_n)=|S/\rho'_n|=n\quad(n\ge1).
\tag{TPS22}
\] The ideals \(P_\tau\) and \(I_0\) have infinite quotient cardinality.
Therefore the original programme ideal zeta and the closed-point residue
Euler product are, for \(\Re s>1\), \[
Z_{\mathrm{ar}}(s)=\sum_{n\ge1}n^{-s}
=\prod_p(1-p^{-s})^{-1}=\zeta(s).
\tag{TPS23}
\] Absolute convergence follows from \(\sum n^{-\sigma}<\infty\) for
\(\sigma>1\), by comparison with \(1+\int_1^\infty x^{-\sigma}dx\).
Finite geometric expansion of a product over primes sums over integers
having only those prime factors. Unique prime factorization then
exhausts the integers; absolute convergence justifies passage to all
primes. The closed points are exactly \(P_p\), with residue
\(\mathbb F_p\), by TPS6 and TPS11. The finite Boolean residue at
\(P_\tau\) contributes no closed-point factor because that point is not
closed.

The cardinality of the actual support-preserving quotient instead is \[
N_{\mathrm{card}}(I_n)=|S/\rho_n|=n+1\quad(n\ge1),
\qquad
Z_{\mathrm{card}}(s)=\sum_{n\ge1}(n+1)^{-s}=\zeta(s)-1.
\tag{TPS24}
\] This weight is not multiplicative, since \((m+1)(n+1)-(mn+1)=m+n\);
even the unit ideal has supported quotient \(G(\mathbb Z/1)=\mathbb B\)
of size two. The mixed-support mechanism at coprime moduli is
TPS20--TPS21.

Define separately the multiplicative extension of the supported
prime-residue sizes: \[
\widetilde N_{\mathrm{sup}}(I_n)=\prod_{p^a\Vert n}(p+1)^a,
\quad \widetilde N_{\mathrm{sup}}(I_1)=1,
\] \[
\boxed{Z_{\mathrm{sup}}(s)=\sum_{n\ge1}\widetilde N_{\mathrm{sup}}(I_n)^{-s}
=\prod_p(1-(p+1)^{-s})^{-1}\quad(\Re s>1).}
\tag{TPS25}
\] This definition agrees with \(|G(\mathbb F_p)|\) on prime ideals. It
is not asserted to equal the cardinality of the composite quotient.
Prime exponents prove multiplicativity of the displayed weight. Since
\(\widetilde N_{\mathrm{sup}}(I_n)\ge n\), its Dirichlet series
converges absolutely in that half-plane. The same finite-prime geometric
expansion as in TPS23 proves its Euler product. Thus TPS22, TPS24 and
TPS25 specify three different counting constructions with their exact
maps and defects retained.

\subsection{7. Complete analytic comparison of the supported
prime-residue
product}\label{complete-analytic-comparison-of-the-supported-prime-residue-product}

On the half-plane \(\Re s>0\), define \[
\ell_p(s)=\log(1-p^{-s})-\log(1-(p+1)^{-s}),
\qquad L_\tau(s)=\sum_p\ell_p(s),\qquad R_\tau(s)=\exp L_\tau(s).
\tag{TPS26}
\] Each logarithm is the analytic branch
\(\log(1-z)=-\sum_{r\ge1}z^r/r\) on \(|z|<1\), so no choice of winding
number is left unspecified. For a real variable \(x\ge2\),
differentiation gives the exact integral \[
\ell_p(s)=-\int_p^{p+1}\frac{s x^{-s-1}}{1-x^{-s}}\,dx.
\tag{TPS27}
\] Fix \(a>0\), put \(d_a=1-2^{-a}>0\), and take \(\Re s\ge a\). Then
\(|1-x^{-s}|\ge d_a\), so \[
|\ell_p(s)|\le\frac{|s|}{d_a}\int_p^{p+1}x^{-a-1}dx,
\qquad
|L_\tau(s)|\le\frac{|s|2^{-a}}{a d_a}.
\tag{TPS28}
\] The prime intervals \([p,p+1]\) have disjoint interiors and are
contained in \([2,\infty)\), which gives the last bound. In particular
the series converges absolutely and uniformly on every compact subset of
the half-plane, and defines a holomorphic function there. Its
exponential is holomorphic and nowhere zero. The product of its
exponential factors is exactly \[
\boxed{R_\tau(s)=\prod_p\frac{1-p^{-s}}{1-(p+1)^{-s}}
\quad(\Re s>0).}
\tag{TPS29}
\] This proves convergence of the ratio product itself beyond the region
where its two separate Euler products converge. For real \(s>0\), each
\(\ell_p(s)<0\), and their sum is finite; hence \(0<R_\tau(s)<1\). Thus
the correction is a nontrivial function while having no zeros in this
domain.

All differentiated prime series converge locally uniformly as well. For
\(j\ge1\), direct differentiation of the logarithmic power series gives
\[
\ell_p^{(j)}(s)=(-1)^{j-1}\sum_{r\ge1}r^{j-1}
\left[(\log p)^j p^{-rs}-(\log(p+1))^j(p+1)^{-rs}\right].
\tag{TPS30}
\] Here is an explicit convergence proof preserving the paired terms. On
a compact set let \(\Re s\ge a>0\), \(|s|\le M\). The derivative with
respect to \(x\) of \((\log x)^j x^{-rs}\) is
\(x^{-rs-1}[j(\log x)^{j-1}-rs(\log x)^j]\). Integrating this derivative
from \(p\) to \(p+1\), and then multiplying by \(r^{j-1}\), bounds the
absolute value of the paired summand by \[
p^{-ra-1}\left[jr^{j-1}(\log(p+1))^{j-1}
+Mr^j(\log(p+1))^j\right].
\] For any fixed nonnegative integer \(d\),
\(\sum_{r\ge1}r^d p^{-ra}\le p^{-a}\sum_{r\ge1}r^d2^{-(r-1)a}\), and the
latter numerical series is finite by the ratio test. Thus the complete
\(p\)-term is bounded by a constant times \(p^{-a-1}(1+(\log(p+1))^j)\).
Summing over all integers bounds the sum over primes, and the integral
test proves convergence. These estimates justify both differentiations
and interchange of the paired prime and prime-power sums on compact
sets.

In particular the exact logarithmic derivative and an explicit
vertical-strip bound are \[
\boxed{r_\tau(s):=\frac{R_\tau'(s)}{R_\tau(s)}
=\sum_p\left[\frac{\log p}{p^s-1}
-\frac{\log(p+1)}{(p+1)^s-1}\right],}
\] \[
r_\tau(s)=\sum_p\int_p^{p+1}x^{-s-1}
\left[\frac{s\log x}{(1-x^{-s})^2}-\frac1{1-x^{-s}}\right]dx,
\] \[
\boxed{|r_\tau(s)|\le
\frac{2^{-a}}{a d_a}
+\frac{|s|2^{-a}}{d_a^2}\left(\frac{\log2}{a}+\frac1{a^2}\right)
\quad(\Re s\ge a>0).}
\tag{TPS31}
\] The integral follows by differentiating \(g_s(x)=\log x/(x^s-1)\) in
\(x\) and writing \(g_s(p)-g_s(p+1)=-\int_p^{p+1}g_s'(x)dx\). Its bound
follows from the disjoint prime intervals, together with
\(\int_2^\infty x^{-a-1}dx=2^{-a}/a\) and
\(\int_2^\infty(\log x)x^{-a-1}dx=2^{-a}(\log2/a+1/a^2)\), the second
obtained by integration by parts. On any fixed strip \(a\le\Re s\le b\),
this is an explicit \(O_{a,b}(1+|\Im s|)\) estimate.

Combining the absolutely convergent Euler products first in \(\Re s>1\)
gives \(Z_{\mathrm{sup}}=\zeta R_\tau\). To supply the continuation used
here explicitly, for \(\Re s>1\) integrate the counting function
\(\lfloor x\rfloor\): \[
\zeta(s)=s\int_1^\infty\lfloor x\rfloor x^{-s-1}dx
=\frac{s}{s-1}-s\int_1^\infty\{x\}x^{-s-1}dx.
\] The first equality follows by interchanging the absolutely convergent
sum of the indicators \(1_{[n,\infty)}\) with the integral. The last
integral and all its derivatives converge uniformly on compact subsets
of \(\Re s>0\), since \(0\le\{x\}<1\) and powers of \(\log x\) are
integrable against \(x^{-a-1}\). Thus it gives a meromorphic
continuation with its sole pole in this domain at \(s=1\), of residue
one. Therefore \[
\boxed{Z_{\mathrm{sup}}(s)=\zeta(s)R_\tau(s)\quad(\Re s>0),
\qquad
-\frac{Z_{\mathrm{sup}}'}{Z_{\mathrm{sup}}}
=-\frac{\zeta'}{\zeta}-r_\tau(s).}
\tag{TPS32}
\] The first equation defines the continuation of the previously defined
prime-residue Euler product. The second holds wherever the denominators
are nonzero and as a meromorphic identity throughout the domain. Since
\(R_\tau\) is holomorphic and nowhere zero, the zeros and poles of
\(Z_{\mathrm{sup}}\) there have exactly the same locations and orders as
those of \(\zeta\); its pole residue at one is \(R_\tau(1)\). The prime
logarithmic weights nevertheless change by the complete function
\(-r_\tau\), with its sign and vertical growth given in TPS31. The
different quotient-cardinality function \(Z_{\mathrm{card}}=\zeta-1\) of
TPS24 has not been substituted for this prime-residue product.

\subsection{8. The full pointed multiplicative
spectrum}\label{the-full-pointed-multiplicative-spectrum}

Let \(M_S\) be the pointed multiplicative monoid underlying \(S\). A
monoid ideal contains \(\tau\) and is closed under multiplication by all
elements; it need not be closed under addition. For a subset \(A\) of
the positive rational primes, define \[
P_A^{\mathrm{mon}}=\{\tau,e\}\cup
\{n^\bullet:n\ne0,\ \exists p\in A\text{ with }p\mid n\}.
\tag{TPS33}
\] The complete monoid-prime spectrum consists of \(P_\tau\) and these
\(P_A^{\mathrm{mon}}\). Indeed a monoid ideal containing any supported
integer also contains \(e\), by multiplication by \(e\). A proper ideal
contains neither unit \(\pm1\). Factor a nonzero integer into its sign
and positive prime powers; repeated primeness shows that it belongs to
the prime exactly when one of its prime factors does. Conversely each
displayed set satisfies primeness by the same factorization, with
products equal to \(e\) and \(\tau\) checked separately. Their inclusion
order is subset inclusion of \(A\), with \(P_\tau\) strictly below every
one. This verifies the existing infinite-dimensional monoid spectrum.

The exact inclusion of the semiring spectrum is \[
X\hookrightarrow\operatorname{Spec}_{\mathrm{mon}}M_S,
\qquad P_\tau\mapsto P_\tau,\quad
P_0\mapsto P_\varnothing^{\mathrm{mon}},\quad
P_p\mapsto P_{\{p\}}^{\mathrm{mon}}.
\tag{TPS34}
\] Its topology is the subspace topology because both basic opens test
exclusion of the same element. A monoid prime \(P_A^{\mathrm{mon}}\) is
additively closed exactly when \(|A|\le1\). These cases are TPS6. If
\(p\ne q\) lie in \(A\), choose integers \(a,b\) with \(ap+bq=1\). Both
supported summands belong to the monoid ideal, while their sum
\(1^\bullet\) does not. This constructs the precise additive defect and
the exact subspace selected by retaining addition.

\subsection{9. Contracted multiplicative algebra with every prime
variable
retained}\label{contracted-multiplicative-algebra-with-every-prime-variable-retained}

Write \([n]\) for the basis element belonging to \(n^\bullet\), and
define \[
\Gamma=\Gamma(S)=\mathbb Z[M_S]/([\tau]),\qquad
E=[0],\quad F=[1]-E,\quad D_n=[n]-E\ (n\in\mathbb Z).
\tag{TPS35}
\] The bracket map is multiplicative. Its ring addition is formal and is
not assumed to be the original addition of \(S\). Since every product
with \([\tau]\) is \([\tau]\), the ideal it generates is its free
rank-one integer span. Thus \(\Gamma\) is free on the integer basis
\([n]\), with \([m][n]=[mn]\), and \(E\ne0\). Multiplication gives
\(E[n]=E\), \(E^2=E\), \(EF=0\), \(F^2=F\), and \(E+F=1\). Also
\(D_0=0\), \(D_mD_n=D_{mn}\), and \(D_1=F\). Comparing coefficients of
the nonzero integer basis elements shows that the \(D_n\), \(n\ne0\),
form a free basis of \(F\Gamma\). Hence \[
\boxed{\Gamma\xrightarrow{\sim}\mathbb Z\times T,\qquad
E\mapsto(1,0),\qquad [n]\mapsto(1,u_n)\ (n\ne0),}
\] \[
T=\mathbb Z[\mathbb Z\setminus\{0\},\times]
=\mathbb Z[\epsilon,X_p:p\text{ prime}]/(\epsilon^2-1),
\quad u_n=\epsilon^{1_{n<0}}\prod_pX_p^{v_p(|n|)}.
\tag{TPS36}
\] The nonzero-integer prime factorization proves the presentation of
\(T\). For a completely explicit inverse, send
\((a,\sum_{n\ne0}b_nu_n)\) to \((a-\sum b_n)E+\sum b_n[n]\). This is
inverse to \(\sum a_n[n]\mapsto(\sum a_n,\sum_{n\ne0}a_nu_n)\), and
multiplicativity follows on basis products. This recovers the previously
proved splitting with its actual labels. In particular \(\epsilon\) is
the sign variable and \(X_p\) is the labelled multiplicative prime, not
the coefficient integer \(p\).

Every ring prime of \(\Gamma\) is exactly one of \[
\mathfrak l\times T\quad(\mathfrak l\in\operatorname{Spec}\mathbb Z),
\qquad
\mathbb Z\times\mathfrak q\quad(\mathfrak q\in\operatorname{Spec}T).
\tag{TPS37}
\] A prime contains at least one of the idempotents \(E,F\), because
their product is zero, and cannot contain both because their sum is one.
Quotienting by the contained idempotent proves the classification and
its topology: the two classes are disjoint open-and-closed components.
Their local rings are \(\mathbb Z_{\mathfrak l}\) and
\(T_{\mathfrak q}\), respectively, since the other idempotent becomes a
unit and kills its complementary factor. The residue fields are the
fraction fields of \(\mathbb Z/\mathfrak l\) and \(T/\mathfrak q\).

Set \(R_0=\mathbb Z[X_p:p\text{ prime}]\). All primes of \(T\) have the
form \[
\mathfrak q_{\sigma,\mathfrak r}=(\epsilon-\sigma)+\mathfrak rT,
\qquad \sigma\in\{1,-1\},\quad\mathfrak r\in\operatorname{Spec}R_0.
\tag{TPS38}
\] Indeed \((\epsilon-1)(\epsilon+1)=0\), so one factor is in the prime,
and quotienting by it gives exactly \(R_0\). The two sign choices give
the same prime exactly when the same \(\mathfrak r\) contains \(2\). The
two sign components thus meet along
\(\operatorname{Spec}\mathbb F_2[X_p]\). The exact ring map \[
T\xrightarrow{\sim}R_0\times_{R_0/2R_0}R_0,
\qquad a+b\epsilon\mapsto(a+b,a-b)
\tag{TPS39}
\] is injective because \(R_0\) has no two-torsion. A pair of
polynomials with equal reductions modulo two has inverse \(a=(f+g)/2\),
\(b=(f-g)/2\), proving surjectivity. These formulas describe all three
irreducible components of \(\operatorname{Spec}\Gamma\): its separate
support copy of \(\operatorname{Spec}\mathbb Z\), and the two sign
copies of \(\operatorname{Spec}R_0\).

Coefficient reduction also gives exact maps retaining the sign
collision: \[
\Gamma\otimes\mathbb F_2\cong
\mathbb F_2\times\mathbb F_2[h,X_p:p\text{ prime}]/(h^2),\qquad h=\epsilon-1,
\] \[
\Gamma\otimes\mathbb F_\ell\cong
\mathbb F_\ell\times\mathbb F_\ell[X_p]\times\mathbb F_\ell[X_p]
\quad(\ell\text{ odd prime}).
\tag{TPS40}
\] In characteristic two, \(\epsilon^2-1=(\epsilon-1)^2\), and the
polynomial basis makes \(h\ne0\), \(h^2=0\). For odd \(\ell\), the two
factors \(\epsilon-1,\epsilon+1\) are comaximal because their difference
is the unit \(2\); evaluation at the two signs and the interpolation
inverse prove the displayed product. The nilpotent here is the retained
characteristic-two sign direction of this monoid algebra.

Contraction of a ring prime through the bracket map always gives a
pointed monoid prime: \[
c:\operatorname{Spec}\Gamma\to\operatorname{Spec}_{\mathrm{mon}}M_S,
\quad c(\mathfrak Q)=\{x\in S:[x]\in\mathfrak Q\}.
\] \[
\boxed{c(\mathfrak l\times T)=P_\tau,\qquad
c(\mathbb Z\times\mathfrak q)=P_{A(\mathfrak q)}^{\mathrm{mon}},
\quad A(\mathfrak q)=\{p:X_p\in\mathfrak q\}.}
\tag{TPS41}
\] For the first formula each supported basis element has first
coordinate one. For the second, \(E\) has second coordinate zero,
\(\epsilon\) is a unit, and primeness tests the individual prime factors
of \(u_n\). Continuity follows because the inverse image of \(D(x)\) is
\(D([x])\). The fibre over \(P_\tau\) is exactly the support component;
the fibre over \(P_A^{\mathrm{mon}}\) is exactly \[
\operatorname{Spec}\bigl(\mathbb Z[\epsilon,X_p^{\pm1}:p\notin A]/(\epsilon^2-1)\bigr).
\tag{TPS42}
\] To prove the latter description, quotient by all \(X_p\) for
\(p\in A\), and localize at all remaining variables. The ordinary prime
correspondence for localization follows directly from taking inverse
images and extending ideals, so this gives precisely primes having the
specified zero-label set. Every set \(A\) occurs: the ideal
\((\epsilon-1,X_p:p\in A)\) has domain quotient
\(\mathbb Z[X_p:p\notin A]\). Adding any coefficient prime \(\ell\)
gives the corresponding example in characteristic \(\ell\).

By TPS34, this contraction is a semiring prime exactly when
\(|A(\mathfrak q)|\le1\), or when it is the support component. In
particular coefficient characteristic and integer prime label are
independent before addition is imposed. The prime
\((2,\epsilon-1)\subset T\) contracts to \(P_0\), since its quotient
\(\mathbb F_2[X_p]\) still has \(X_2\ne0\). The homomorphism to
\(\mathbb F_2\) sending \(\epsilon\mapsto1\), \(X_2,X_3\mapsto0\), and
every other prime variable to one has prime kernel contracting to
\(P_{\{2,3\}}^{\mathrm{mon}}\). That set contains \(3^\bullet\) and
\((-2)^\bullet\), but not their sum \(1^\bullet\). Thus the failure of a
map from the entire ring spectrum to the semiring spectrum is witnessed
by these exact additive relations, and TPS34 identifies precisely the
locus where its set-theoretic contraction is an additive prime.

\subsection{10. The chosen addition relations and all their induced
spectrum
maps}\label{the-chosen-addition-relations-and-all-their-induced-spectrum-maps}

Define evaluation and its exact kernel by \[
\eta:T\twoheadrightarrow\mathbb Z,\quad
\epsilon\mapsto-1,\quad X_p\mapsto p,\qquad
K=(\epsilon+1,X_p-p:p\text{ prime}).
\tag{TPS43}
\] The quotient by the displayed generators makes every variable its
indicated constant, and thus has mutually inverse maps with
\(\mathbb Z\); this proves the equality of the kernel. Every polynomial
has finitely many variables, so the argument covers the entire
infinitely generated polynomial ring.

Let \(J_{\mathrm{ar}}\) be generated in \(\Gamma\) by the original
additive relations \([m+n]-[m]-[n]\), and let \(J_{\mathrm{sup}}\) be
generated by \([m+n]-[m]-[n]+E\), for all integers \(m,n\). Then \[
\boxed{J_{\mathrm{sup}}=0\times K,\quad
J_{\mathrm{ar}}=\mathbb Z\times K,\quad
\Gamma/J_{\mathrm{sup}}\cong\mathbb Z\times\mathbb Z,\quad
\Gamma/J_{\mathrm{ar}}\cong\mathbb Z.}
\tag{TPS44}
\] For a full proof, the supported relation is \(D_{m+n}-D_m-D_n\), and
has coordinates \((0,\text{an element of }K)\), proving one containment.
Telescoping the relations for \((j,1)\), \(1\le j<n\), gives \(D_n-nF\)
for every positive \(n\); the relation for \((-1,1)\) gives
\(D_{-1}+F\). These include the lifts of every generator \(X_p-p\) and
\(\epsilon+1\), proving the other containment. The relation for
\((0,0)\) in the original-addition ideal equals \(-E\). Thus
\(J_{\mathrm{ar}}=J_{\mathrm{sup}}+(E)\), proving all of TPS44. In the
original brackets the generating sets can be written \[
J_{\mathrm{sup}}=([-1]+[1]-2E,\ [p]-p[1]+(p-1)E:p\text{ prime}),
\] \[
J_{\mathrm{ar}}=(E,\ [-1]+[1],\ [p]-p[1]:p\text{ prime}).
\tag{TPS45}
\] These exhibit exactly which prime-label variables become their
numerical integer values.

On spectra, the arithmetic quotient is the closed immersion with image
\[
\{\mathbb Z\times K\}\cup
\{\mathbb Z\times(K+(p)):p\text{ prime}\};
\] its contractions to \(S\) are \(P_0,P_p\). The supported-addition
quotient retains this closed copy and the entire separate support
component \(\{\mathfrak l\times T\}\), all of which contracts to
\(P_\tau\). These statements follow by the elementary prime
correspondence for a quotient: primes in the quotient correspond exactly
to primes containing its kernel. They preserve the coefficient primes on
the support component instead of identifying them with integer-labelled
closed points of \(S\).

The maps themselves are \[
\Gamma\to\mathbb Z\times\mathbb Z,\quad (a,t)\mapsto(a,\eta(t)),\qquad
[n]\mapsto(1,n),\quad[\tau]\mapsto(0,0),
\] \[
\Gamma\to\mathbb Z,\quad(a,t)\mapsto\eta(t),\qquad
[n]\mapsto n,\quad[\tau]\mapsto0,
\] \[
\Gamma\to\mathbb Z,\quad(a,t)\mapsto a,\qquad
[n]\mapsto1,\quad[\tau]\mapsto0.
\tag{TPS46}
\] The last is the ring extension of multiplicative Boolean support; it
descends through \(J_{\mathrm{sup}}\), while an original additive
relation has image \(-1\) under it. The first map is injective on the
original set \(S\) and preserves multiplication, but its exact additive
defect is \[
i(x+y)=i(x)+i(y)-\chi(x)\chi(y)(1,0).
\tag{TPS47}
\] If both arguments are supported this subtracts exactly one support
coordinate; if either is \(\tau\) it subtracts zero. This proves all
cases. Replacing the first coefficient ring by \(\mathbb B\) instead
gives the actual semiring embedding of TPS3, with its Boolean addition.
Conversely every unital semiring map from \(S\) into a ring sends \(e\)
to zero, since \(e+e=e\) and additive cancellation hold in the target.
It therefore factors through \(\pi\). This proves exactly why the
original-addition ring quotient removes the supported idempotent and why
the supported-addition quotient records a different displayed relation.

\subsection{11. Applying the multiplicative algebra to the supported
finite
residue}\label{applying-the-multiplicative-algebra-to-the-supported-finite-residue}

The support-preserving semiring quotient \(q_p^{\mathrm{sup}}\) has its
own contracted multiplicative algebra: \[
\boxed{\Gamma(G(\mathbb F_p))\cong
\mathbb Z\times\mathbb Z[\mathbb F_p^\times].}
\tag{TPS48}
\] No generator of \(\mathbb F_p^\times\) is chosen. Write
\(\langle a\rangle\) for its group-ring basis, and \(E_p=[0^\bullet]\).
The isomorphism sends \(E_p\mapsto(1,0)\) and
\([a^\bullet]\mapsto(1,\langle a\rangle)\) for \(a\ne0\). Its inverse
sends \((z,\sum c_a\langle a\rangle)\) to
\((z-\sum c_a)E_p+\sum c_a[a^\bullet]\). Multiplication of basis
elements proves it is an isomorphism exactly as in TPS36.

The induced map is
\(\Gamma(q_p^{\mathrm{sup}})=\mathrm{id}_{\mathbb Z}\times\psi_p\),
where \[
\psi_p(u_n)=\begin{cases}0,&p\mid n,\\\langle\bar n\rangle,&p\nmid n,\end{cases}
\quad
\psi_p(X_p)=0,\quad\psi_p(X_\ell)=\langle\bar\ell\rangle\ (\ell\ne p),\quad
\psi_p(\epsilon)=\langle-1\rangle.
\tag{TPS49}
\] This is multiplicative because reduction modulo \(p\) preserves
products and a product is zero precisely when a factor is zero. It is
onto since every nonzero residue has an integer representative. Its
exact kernel is \[
L_p=(X_p,\ u_m-u_n:
m,n\ne0,\ p\nmid mn,\ m\equiv n\pmod p)\subset T,
\qquad\ker\Gamma(q_p^{\mathrm{sup}})=0\times L_p.
\tag{TPS50}
\] All generators map to zero. Conversely choose representatives
\(r_a\in\{1,\ldots,p-1\}\) of the nonzero residues. Modulo the displayed
ideal, each basis \(u_n\) is zero when \(p\mid n\), since it is
divisible by \(X_p\), and otherwise equals \(u_{r_a}\). These surviving
representatives map to the free group-ring basis, so any element mapping
to zero has every coefficient zero. This proves the claimed kernel. In
particular \(p\notin L_p\): the target group ring is a free abelian
group and has characteristic zero, even though its basis labels come
from a field of characteristic \(p\).

Arithmetic evaluation is the further surjection \[
\rho_p:\mathbb Z[\mathbb F_p^\times]\twoheadrightarrow\mathbb F_p,
\quad\langle a\rangle\mapsto a,
\qquad
H_p=(p\langle1\rangle,\ \langle a\rangle-r_a\langle1\rangle:a\ne0).
\tag{TPS51}
\] The displayed generators give its exact kernel because their quotient
makes every basis vector the stated scalar and imposes precisely
coefficient characteristic \(p\). Changing representatives alters them
only by multiples of \(p\langle1\rangle\). Hence \[
\psi_p^{-1}(H_p)=K+(p),\qquad L_p\subseteq K+(p),
\] \[
\mathbb Z\times T\twoheadrightarrow
\mathbb Z\times\mathbb Z[\mathbb F_p^\times]\twoheadrightarrow
\mathbb Z\times\mathbb F_p
\tag{TPS52}
\] has composite kernel \(0\times(K+(p))\). To identify the latter map
with its addition relations, write \(d_a=[a^\bullet]-E_p\). The
supported relations are \(d_{a+b}=d_a+d_b\); repeated addition gives
\(d_{r\bmod p}=r d_1\), including \(p d_1=d_0=0\). They therefore
generate exactly \(0\times H_p\). The original addition relations also
kill \(E_p\) and yield \(\mathbb F_p\). This distinguishes the finite
supported semiring, its characteristic-zero multiplicative algebra, and
its supported or ordinary additive ring quotients by exact surjections
and kernels.

\subsection{12. Integral and adelic coordinate primes with their exact
contractions}\label{integral-and-adelic-coordinate-primes-with-their-exact-contractions}

The following algebraic maps specify the coordinate-prime comparison;
they do not identify an adelic orbit space with the scalar prime
spectrum. Define \[
\mathbb Z_p=\varprojlim_j\mathbb Z/p^j\mathbb Z,\qquad
\mathbb Q_p=\mathbb Z_p[1/p],\qquad
\widehat{\mathbb Z}=\prod_p\mathbb Z_p,
\] \[
\mathbb A_f=\{(a_p)\in\prod_p\mathbb Q_p:a_p\in\mathbb Z_p
\text{ for all but finitely many }p\},\qquad
\mathbb A_{\mathbb Q}=\mathbb R\times\mathbb A_f.
\tag{TPS53}
\] For the algebraic facts needed here, a nonzero element of
\(\mathbb Z_p\) has a unique form \(p^k u\) with \(u\) a unit: choose
its first nonzero residue digit, divide its compatible representatives
by \(p^k\), and invert the resulting nonzero residues modulo each
\(p^j\). Their inverses are compatible, giving the inverse in
\(\mathbb Z_p\). Multiplying these expressions shows that valuations
add, so \(\mathbb Z_p\) is a domain and \(\mathbb Q_p\) is its fraction
field. Its residue modulo \(p\) is \(\mathbb F_p\). The diagonal map
\(\mathbb Z\to\mathbb Z_p\) is injective, since an integer divisible by
every \(p^j\) is zero. The product description of
\(\widehat{\mathbb Z}\) agrees with its inverse-limit-over-all-moduli
description: finite Chinese remainder isomorphisms, proved as in TPS20,
reconstruct each compatible modulus from the prime-power coordinates.

Every unital ring homomorphism \(f:A\to B\) induces a unital semiring
homomorphism \(G(f):G(A)\to G(B)\), sending \(\tau\) to \(\tau\) and
\(a^\bullet\) to \(f(a)^\bullet\). The defining laws TPS2 prove this on
every pair of elements. In particular the canonical diagrams are \[
S\hookrightarrow G(\widehat{\mathbb Z})\hookrightarrow G(\mathbb A_f),
\qquad
\begin{array}{ccc}
S&\longrightarrow&G(\mathbb A_{\mathbb Q})\\
\downarrow&&\downarrow\scriptstyle G(\mathrm{pr}_f)\\
G(\widehat{\mathbb Z})&\longrightarrow&G(\mathbb A_f).
\end{array}
\tag{TPS54}
\] Denote the diagonal map \(S\to G(\widehat{\mathbb Z})\) by
\(\iota_{\widehat{\mathbb Z}}\). The source maps are the diagonal
integer maps, and the last horizontal map is inclusion of integral
finite adeles. All are unital. The square commutes on every supported
integer and on \(\tau\), so everywhere. The full adele comparison uses
this square; the rule \(a_f\mapsto(0,a_f)\) would have unit image
\((0,1)\), so it is not a unital replacement for a map into the full
adeles.

For a fixed rational prime \(p\), define two integral coordinate ideals
\[
\mathfrak q_p=\ker(\widehat{\mathbb Z}\to\mathbb Z_p),\qquad
\mathfrak m_p=\{a\in\widehat{\mathbb Z}:a_p\in p\mathbb Z_p\}.
\tag{TPS55}
\] Coordinate projection is onto, so their quotients are \(\mathbb Z_p\)
and \(\mathbb F_p\). Thus \(\mathfrak q_p\subsetneq\mathfrak m_p\), the
first is prime and the second maximal. Strictness is witnessed by the
element with coordinate \(p\) at \(p\) and zero at all other places. For
any ring \(A\), the ideals of \(G(A)\) are \(\{\tau\}\) and
\(\{\tau\}\cup J^\bullet\) for ring ideals \(J\); the proof of TPS5
works unchanged until the integer principal-ideal classification, which
is unnecessary here. Such a lifted proper ideal is prime exactly when
\(J\) is prime, by testing supported products. Consequently the exact
contractions along TPS54 are \[
\boxed{\iota_{\widehat{\mathbb Z}}^{-1}(\{\tau\}\cup\mathfrak q_p^\bullet)=P_0,
\qquad
\iota_{\widehat{\mathbb Z}}^{-1}(\{\tau\}\cup\mathfrak m_p^\bullet)=P_p.}
\tag{TPS56}
\] The first equality follows from injectivity of
\(\mathbb Z\to\mathbb Z_p\); the second from the condition \(p\mid n\)
on an integer. The support prime \(\{\tau\}\) itself contracts to
\(P_\tau\). Thus coordinate supported zero and the external unsupported
element give different prime contractions.

There is a useful exact localization of the entire integral coordinate
ring: \[
\widehat{\mathbb Z}[1/n:n\in\mathbb Z\setminus\{0\}]
\xrightarrow{\sim}\mathbb A_f.
\tag{TPS57}
\] Each rational integer denominator is invertible in \(\mathbb A_f\),
and its reciprocal is integral outside its finitely many prime divisors.
Conversely a finite adele has finitely many negative valuations;
multiply by the product of the corresponding prime powers to make all
its coordinates integral. It is therefore a fraction from the left side.
Multiplication by a nonzero diagonal integer is injective in
\(\widehat{\mathbb Z}\), since each coordinate ring is a domain. This
proves injectivity of the fraction map as well as surjectivity.

Let \(e_p\in\widehat{\mathbb Z}\) be the idempotent which is one at
\(p\) and zero elsewhere. Then
\(\mathfrak q_p=(1-e_p)\widehat{\mathbb Z}\), and extension through
TPS57 gives \[
\mathfrak q_p\mathbb A_f=(1-e_p)\mathbb A_f
=\ker(\mathbb A_f\to\mathbb Q_p),\qquad
\mathfrak m_p\mathbb A_f=\mathbb A_f.
\tag{TPS58}
\] The first equality is evaluation at the \(p\)-coordinate; every adele
with that coordinate zero is fixed by multiplication by \(1-e_p\). The
second holds because the diagonal integer \(p\) belongs to
\(\mathfrak m_p\) and is a unit in \(\mathbb A_f\). Thus the maximal
integral residue ideal defining \(\mathbb F_p\) does not extend to a
proper finite-adele ideal, whereas the zero-coordinate ideal survives
and has quotient \(\mathbb Q_p\).

At each finite or infinite place \(v\), let \(K_v=\mathbb Q_p\) or
\(\mathbb R\), and set
\(\mathfrak n_v=\ker(\mathbb A_{\mathbb Q}\to K_v)\). Each coordinate
projection is onto, so \(\mathfrak n_v\) is maximal with quotient
\(K_v\). The diagonal integers inject into every \(K_v\), proving \[
\boxed{\bigl(S\to G(\mathbb A_{\mathbb Q})\bigr)^{-1}
(\{\tau\}\cup\mathfrak n_v^\bullet)=P_0\quad\text{for every place }v.}
\tag{TPS59}
\] These coordinate ideals therefore retain their distinct place labels
in the adelic ring while all contracting to the same scalar
arithmetic-generic prime. This is an explicit many-to-one contraction,
not an identification of a place-orbit construction with the
rational-prime closed points. Only the specified coordinate ideals are
classified in TPS55--TPS59; no classification of all primes of an
infinite product or restricted product is asserted.

\subsection{13. Reading scope and exact return to the original
programme}\label{reading-scope-and-exact-return-to-the-original-programme}

TPS1--TPS18 verify the existing original semiring, spectrum, topology,
congruences and stalks with their two coordinates. TPS19--TPS25 retain
the mixed-support sectors and distinguish the three actual counting
constructions. TPS26--TPS32 prove the analytic ratio, its complete
differentiated series and explicit contour-relevant growth. TPS33--TPS52
retain all prime variables and compute the multiplicative/additive
algebra maps, their kernels and their spectrum contractions.
TPS53--TPS59 give the integral and adelic coordinate maps needed to
avoid identifying supported coordinate zero with \(\tau\), or
identifying the place-labelled coordinate ideals with their scalar
contractions.

The original ideal and stalk theorems and the contracted-ring splitting
are prior programme mathematics with the source locators at the start of
this note. The complete finite and analytic proofs above are the present
verification and propagation. The note makes no claim that a spectral
interpretation of zeta zeros, an adelic trace formula, or the Riemann
hypothesis follows from these prime classifications alone.

\end{document}
